\documentclass[12pt,a4paper]{article}
\usepackage[utf8]{inputenc}
\usepackage[vietnamese]{babel}
\usepackage{amsmath,amsfonts,amssymb}
\usepackage{graphicx}
\usepackage{algorithm}
\usepackage{algorithmic}
\usepackage{booktabs}
\usepackage{multirow}
\usepackage{geometry}
\usepackage{hyperref}
\usepackage{listings}
\usepackage{xcolor}
\usepackage{float}

\geometry{margin=2.5cm}

\lstset{
    basicstyle=\ttfamily\small,
    keywordstyle=\color{blue},
    commentstyle=\color{green!60!black},
    stringstyle=\color{red},
    numbers=left,
    numberstyle=\tiny,
    breaklines=true,
    frame=single
}

\title{\textbf{BÁO CÁO KỸ THUẬT} \\ \Large Giải Bài Toán Pickup-Delivery với Drone Resupply \\ sử dụng Thuật toán Di truyền kết hợp Tabu Search}
\author{Drone Resupply PDP Solver}
\date{\today}

\begin{document}

\maketitle

\begin{abstract}
Báo cáo này trình bày chi tiết về hệ thống giải bài toán Pickup-Delivery Problem (PDP) với Drone Resupply sử dụng thuật toán lai ghép giữa Genetic Algorithm (GA) và Tabu Search. Chương trình áp dụng mô hình One-Rendezvous cho phép drone tiếp tế hàng hóa cho xe tải tại các điểm gặp nhau, kết hợp với các kỹ thuật consolidation thông minh để gom nhiều khách hàng trong một chuyến bay drone. Kết quả thực nghiệm trên các bộ dữ liệu benchmark cho thấy thuật toán đạt hiệu quả cao với tất cả các solution đều valid và thời gian hoàn thành (C\_max) được tối ưu.
\end{abstract}

\tableofcontents
\newpage

%==============================================================================
\section{Giới thiệu}
%==============================================================================

\subsection{Bối cảnh và động lực}

Trong bối cảnh logistics hiện đại, việc kết hợp xe tải truyền thống với drone đang trở thành xu hướng quan trọng để tối ưu hóa việc giao hàng chặng cuối (last-mile delivery). Drone có ưu điểm về tốc độ và khả năng bay theo đường thẳng (Euclidean), trong khi xe tải có sức chứa lớn hơn và phù hợp với các tuyến đường đô thị (Manhattan).

Bài toán Pickup-Delivery Problem with Drone Resupply (PDP-DR) là một biến thể phức tạp của bài toán định tuyến xe cổ điển, trong đó:
\begin{itemize}
    \item Có các cặp pickup-delivery cần được phục vụ theo thứ tự
    \item Drone có thể bay từ depot để tiếp tế hàng cho xe tải tại các điểm hẹn
    \item Mục tiêu là tối thiểu hóa thời gian hoàn thành tổng thể (makespan - $C_{max}$)
\end{itemize}

\subsection{Mục tiêu của chương trình}

Chương trình được phát triển nhằm:
\begin{enumerate}
    \item Giải quyết bài toán PDP-DR với mô hình One-Rendezvous
    \item Áp dụng thuật toán metaheuristic lai ghép GA + Tabu Search
    \item Hỗ trợ consolidation - gom nhiều khách hàng trong một chuyến drone
    \item Đảm bảo tính khả thi của solution (ràng buộc thời gian, capacity, endurance)
\end{enumerate}

%==============================================================================
\section{Mô hình Bài toán}
%==============================================================================

\subsection{Định nghĩa bài toán}

Bài toán xét một đội xe tải (truck) và drone xuất phát từ depot $0$ để phục vụ tập
khách hàng $N$. Mỗi khách hàng được xe tải ghé đúng một lần. Với khách hàng giao hàng
trực tiếp $i \in N^D$ (loại \texttt{D}), kiện hàng của $i$ có thể được chất lên xe tải
tại depot đầu một chuyến truck, hoặc được drone mang từ depot tới một điểm hẹn trên
hành trình truck. Kiện loại \texttt{D} (gọi tắt C1) phải được giao trước khi truck
quay lại depot kết thúc chuyến đã nhận kiện đó. Với cặp pickup--delivery
$(p,\ell)\in\mathcal A$ (loại \texttt{P}/\texttt{DL}, gọi tắt C2), pickup $p$ và
delivery $\ell$ phải nằm trên cùng một truck, và delivery phải nằm sau pickup trong
toàn hành trình của truck đó.

Mỗi truck có thể đi nhiều chuyến, mỗi chuyến có dạng $0 \to \cdots \to 0$. Khi truck
quay lại depot để bắt đầu chuyến tiếp theo, truck không bắt buộc phải rỗng: hàng đã
pickup ở một khách hàng (C2) có thể được giữ trên truck để giao cho delivery tương ứng
ở chuyến sau, miễn là capacity truck luôn được tôn trọng; hàng C1 thì không được phép
đi qua một lần quay về depot như vậy. Lần xuất phát đầu tiên từ depot không tính thời
gian bốc dỡ truck; từ chuyến thứ hai trở đi có tính thời gian bốc dỡ tại depot
$\delta_T$. Lần cuối truck quay về depot để kết thúc hành trình không tính thêm thời
gian bốc dỡ.

Mỗi chuyến drone chỉ có \emph{một} điểm hẹn với một truck tại một khách hàng (mô hình
\emph{one-rendezvous}): drone xuất phát từ depot, bay tới điểm hẹn $h$, bàn giao hàng
cho truck, rồi quay về depot ($0 \to h \to 0$); mô hình không dùng multi-rendezvous
cho drone (drone không tiếp tục bay tới một truck/điểm hẹn thứ hai trong cùng chuyến
bay). Đây là điểm khác biệt chính so với mô hình VRP-MRDR (multi-rendezvous) được đề
xuất gần đây bởi Tran et al. (2026) -- xem mục~\ref{sec:related-work}.

\subsection{Công trình liên quan}
\label{sec:related-work}

Bài toán truck--drone resupply với release date (khác với drone giao hàng trực tiếp)
được giới thiệu bởi Pina-Pardo, Silva \& Smith (2021), với 1 truck, 1 drone,
uncapacitated truck, mục tiêu minimize makespan, giải bằng MILP kết hợp decomposition.
Pina-Pardo et al. (2024a) mở rộng sang nhiều drone (dynamic). Pina-Pardo et al.
(2024b) mở rộng tiếp sang nhiều truck và nhiều drone có capacity (\emph{Fleet resupply
by drones for last-mile delivery}), vẫn trong mô hình \emph{một rendezvous} mỗi chuyến
bay -- đây là mô hình gần nhất với bài toán được xét trong báo cáo này.

Gần đây, Tran, Nguyen, Binh, Tran \& Nguyen (2026, \emph{Transportation Research Part
E}) đề xuất VRP-MRDR (\emph{Vehicle Routing Problem with Release Dates and
Multi-Rendezvous Drone Resupply}), cho phép một drone tiếp tế tuần tự \textbf{nhiều}
truck trong cùng một chuyến bay, cùng với một MILP đầy đủ và một framework Adaptive
Tabu Search (ATS) hai lớp (routing + loading subproblem bằng CP-SAT). Kiến trúc
TS-ALNS-LSP của mục~\ref{sec:algorithm} được xây dựng theo đúng tinh thần của ATS
trong bài báo này, nhưng hiện mới chỉ hỗ trợ one-rendezvous (mục~2); mở rộng sang
multi-rendezvous là một hướng phát triển trực tiếp (xem mục~\ref{sec:future-work}).

\paragraph{Khoảng cách và vận tốc.} Truck dùng khoảng cách Manhattan, vận tốc $30$
km/h. Drone dùng khoảng cách Euclidean, vận tốc $60$ km/h. Thời gian di chuyển tính
theo phút:
\[
  t^T_{ij}=\frac{|x_i-x_j|+|y_i-y_j|}{30}\cdot 60,\qquad
  t^D_{ij}=\frac{\sqrt{(x_i-x_j)^2+(y_i-y_j)^2}}{60}\cdot 60.
\]

\subsection{Các thông số đầu vào}

\begin{table}[H]
\centering
\caption{Các thông số của bài toán}
\begin{tabular}{lll}
\toprule
\textbf{Ký hiệu} & \textbf{Mô tả} & \textbf{Giá trị mặc định} \\
\midrule
$K$ & Số lượng xe tải & 2 \\
$D$ & Số lượng drone & 2 \\
$Q_t$ & Capacity của xe tải & 50 \\
$Q_d$ & Capacity của drone & 2 (instance nhỏ) / 10 (instance lớn) \\
$v_t$ & Tốc độ xe tải & 30 km/h \\
$v_d$ & Tốc độ drone & 60 km/h \\
$E$ & Drone endurance & 90 phút \\
$s_t$ & Truck service time & 3 phút \\
$s_r$ & Resupply time & 5 phút \\
$s_{depot}$ & Depot receive time & 5 phút \\
$s_{load}$ & Drone load time tại depot & 5 phút \\
$r_i$ & Ready time của khách hàng $i$ & Từ dữ liệu \\
\bottomrule
\end{tabular}
\end{table}

\begin{table}[H]
\centering
\caption{Cấu hình các vị trí Depot}
\begin{tabular}{lll}
\toprule
\textbf{Depot} & \textbf{Tọa độ (x, y)} & \textbf{Mô tả} \\
\midrule
Depot Center & (10.0, 10.0) & Depot trung tâm (mặc định) \\
Depot Border & (0.0, 10.0) & Depot biên \\
Depot Outside & (-10.0, 10.0) & Depot ngoài khu vực \\
\bottomrule
\end{tabular}
\end{table}

\subsection{Tập chỉ số và tham số}

\begin{table}[H]
\centering
\begin{tabular}{ll}
\toprule
Ký hiệu & Ý nghĩa\\
\midrule
$0$ & depot\\
$N$ & tập khách hàng\\
$N^D$ & tập khách hàng giao hàng trực tiếp, loại \texttt{D} (C1)\\
$N^P$ & tập node pickup, loại \texttt{P}\\
$N^{DL}$ & tập node delivery của pickup--delivery, loại \texttt{DL}\\
$\mathcal A$ & tập cặp pickup--delivery $(p,\ell)$ (C2)\\
$K$ & tập truck\\
$R$ & tập drone\\
$H=\{1,\ldots,H_{\max}\}$ & tập chuyến tiềm năng của mỗi truck\\
$P=\{1,\ldots,P_{\max}\}$ & tập vị trí tiềm năng trong một chuyến truck\\
$S=\{1,\ldots,S_{\max}\}$ & tập chuyến bay tiềm năng của mỗi drone\\
\bottomrule
\end{tabular}
\end{table}

\begin{table}[H]
\centering
\begin{tabular}{ll}
\toprule
Ký hiệu & Ý nghĩa\\
\midrule
$q_i$ & khối lượng kiện hàng của khách $i$\\
$r_i$ & release date của kiện hàng $i$\\
$Q_T$ & tải trọng truck\\
$Q_D$ & tải trọng drone\\
$s_T$ & thời gian truck phục vụ một khách hàng\\
$\delta_m$ & thời gian bàn giao drone--truck tại điểm hẹn\\
$\delta_d$ & thời gian drone chất hàng tại depot trước mỗi chuyến bay\\
$\delta_T$ & thời gian truck bốc dỡ tại depot trước chuyến truck thứ hai trở đi\\
$E$ & endurance tối đa của một chuyến bay drone\\
$M$ & hằng số Big-M đủ lớn\\
$t^T_{ij}$ & thời gian truck đi từ $i$ đến $j$\\
$t^D_{ij}$ & thời gian drone bay từ $i$ đến $j$\\
\bottomrule
\end{tabular}
\end{table}

\subsection{Biến quyết định}

\textbf{Route truck.}
\[
y_{kmpi}=\begin{cases}1, & \text{truck }k\text{ phục vụ khách }i\text{ ở vị trí }p\text{ trong chuyến }m,\\0, & \text{ngược lại,}\end{cases}
\qquad
u_{kmp}=\begin{cases}1, & \text{vị trí }p\text{ của chuyến }m\text{ của truck }k\text{ được dùng,}\\0, & \text{ngược lại,}\end{cases}
\]
\[
g_{km}=\begin{cases}1, & \text{chuyến }m\text{ của truck }k\text{ được dùng,}\\0, & \text{ngược lại,}\end{cases}
\qquad
z_{kmpij}=\begin{cases}1, & i\text{ ở vị trí }p\text{ và }j\text{ ở }p+1\text{ trong chuyến }m\text{ của truck }k,\\0, & \text{ngược lại,}\end{cases}
\]
\[
\lambda_{kmpi}=\begin{cases}1, & i\text{ là khách cuối của chuyến }m\text{ của truck }k\text{ tại vị trí }p,\\0, & \text{ngược lại.}\end{cases}
\]

\textbf{Nguồn hàng và drone.}
\[
a_{kmi}=\begin{cases}1, & \text{kiện C1 }i\text{ chất lên truck }k\text{ tại depot đầu chuyến }m,\\0, & \text{ngược lại,}\end{cases}
\qquad
b_{kmpi}=\begin{cases}1, & \text{kiện C1 }i\text{ drone giao cho truck }k\text{ tại vị trí }p\text{ của chuyến }m,\\0, & \text{ngược lại,}\end{cases}
\]
\[
\rho_{kmp}=\begin{cases}1, & \text{có chuyến drone hẹn truck }k\text{ tại vị trí }p\text{ của chuyến }m,\\0, & \text{ngược lại,}\end{cases}
\qquad
w_{dskmpj}=\begin{cases}1, & \text{chuyến bay }s\text{ của drone }d\text{ gặp truck }k\text{ trong chuyến }m\text{ tại vị trí }p\text{, node }j,\\0, & \text{ngược lại,}\end{cases}
\]
\[
\eta_{ds}=\begin{cases}1, & \text{chuyến bay }s\text{ của drone }d\text{ được dùng,}\\0, & \text{ngược lại.}\end{cases}
\]

\textbf{Thời gian và tải.} $T^0_{km}, A_{kmp}, D_{kmp}, B_{km} \ge 0$ lần lượt là
thời điểm truck $k$ rời depot bắt đầu chuyến $m$, thời điểm đến/rời vị trí $p$, và
thời điểm truck quay về depot kết thúc chuyến $m$. $F^{dep}_{ds}, F^{arr}_{ds},
F^{ret}_{ds} \ge 0$ là thời điểm drone $d$ khởi hành/đến điểm hẹn/quay về depot ở
chuyến bay $s$. $L_{kmp}, L^{end}_{km} \ge 0$ là tải trên truck $k$ sau khi rời vị trí
$p$ (và khi kết thúc chuyến $m$) -- tải còn lại cuối chuyến chỉ có thể là hàng pickup
C2 chưa tới delivery tương ứng. $C_{\max} \ge 0$ là makespan.

\subsection{Hàm mục tiêu}

\begin{equation}
  \min C_{\max}.
\end{equation}

\subsection{Các ràng buộc}

\subsubsection{Cấu trúc route truck nhiều chuyến}

\begin{equation}
  \sum_{i\in N} y_{kmpi}=u_{kmp}, \qquad \forall k\in K,\ m\in H,\ p\in P.
\end{equation}
\begin{equation}
  u_{kmp}\le g_{km}, \qquad \forall k\in K,\ m\in H,\ p\in P.
\end{equation}
\begin{equation}
  u_{km,p+1}\le u_{kmp}, \qquad \forall k\in K,\ m\in H,\ p=1,\ldots,P_{\max}-1.
\end{equation}
\begin{equation}
  g_{k,m+1}\le g_{km}, \qquad \forall k\in K,\ m=1,\ldots,H_{\max}-1.
\end{equation}
\begin{equation}
  \sum_{k\in K}\sum_{m\in H}\sum_{p\in P} y_{kmpi}=1, \qquad \forall i\in N.
\end{equation}
\begin{equation}
  z_{kmpij}\le y_{kmpi},\quad z_{kmpij}\le y_{km,p+1,j},\quad z_{kmpij}\ge y_{kmpi}+y_{km,p+1,j}-1,
\end{equation}
\[
  \forall k\in K,\ m\in H,\ p=1,\ldots,P_{\max}-1,\ i,j\in N.
\]
\begin{equation}
  \lambda_{kmpi}\le y_{kmpi},\quad \lambda_{kmpi}\le u_{kmp}-u_{km,p+1},\quad \lambda_{kmpi}\ge y_{kmpi}+u_{kmp}-u_{km,p+1}-1,
\end{equation}
\[
  \forall k\in K,\ m\in H,\ p=1,\ldots,P_{\max}-1,\ i\in N.
\]
\begin{equation}
  \lambda_{km,P_{\max},i}=y_{km,P_{\max},i}, \qquad \forall k\in K,\ m\in H,\ i\in N.
\end{equation}

\subsubsection{Pickup--delivery (C2)}

Pickup và delivery của cùng một cặp phải nằm trên cùng một truck, nhưng có thể nằm
ở hai chuyến khác nhau:
\begin{equation}
  \sum_{m\in H}\sum_{p\in P}y_{kmp p'} = \sum_{m\in H}\sum_{p\in P}y_{kmp \ell},
  \qquad \forall k\in K,\ (p',\ell)\in\mathcal A.
\end{equation}
Pickup phải đứng trước delivery trong toàn hành trình của truck. Thứ tự toàn hành
trình của vị trí $(m,p)$ được tính bởi $(m-1)P_{\max}+p$:
\begin{equation}
\begin{aligned}
  &\sum_{m\in H}\sum_{p\in P}\big((m-1)P_{\max}+p\big)y_{kmp p'}+1
  \le
  \sum_{m\in H}\sum_{p\in P}\big((m-1)P_{\max}+p\big)y_{kmp\ell}\\
  &\quad +M\left(1-\sum_{m\in H}\sum_{p\in P}y_{kmp p'}\right),
  \qquad \forall k\in K,\ (p',\ell)\in\mathcal A.
\end{aligned}
\end{equation}

\subsubsection{Nguồn hàng cho khách C1}

\begin{equation}
  a_{kmi}+\sum_{p\in P}b_{kmpi} = \sum_{p\in P}y_{kmpi}, \qquad \forall k\in K,\ m\in H,\ i\in N^D.
\end{equation}
Do cả nguồn hàng và lần phục vụ đều mang cùng chỉ số chuyến $m$, phương trình trên
buộc mọi kiện C1 được chất tại depot hoặc nhận từ drone phải được giao trong chính
chuyến truck đó; kiện C1 không thể được mang qua một lần quay về depot.
\begin{equation}
  b_{kmpi}\le \sum_{\substack{q\in P\\q\ge p}} y_{kmqi}, \qquad \forall k\in K,\ m\in H,\ p\in P,\ i\in N^D.
\end{equation}
\begin{equation}
  b_{kmpi}\le \rho_{kmp}, \qquad \forall k\in K,\ m\in H,\ p\in P,\ i\in N^D.
\end{equation}
\begin{equation}
  \rho_{kmp}\le u_{kmp}, \qquad \forall k\in K,\ m\in H,\ p\in P.
\end{equation}

\subsubsection{Drone capacity và gán chuyến drone}

\begin{equation}
  \sum_{i\in N^D}q_i b_{kmpi}\le Q_D\rho_{kmp}, \qquad \forall k\in K,\ m\in H,\ p\in P.
\end{equation}
\begin{equation}
  \sum_{d\in R}\sum_{s\in S}\sum_{j\in N} w_{dskmpj} = \rho_{kmp}, \qquad \forall k\in K,\ m\in H,\ p\in P.
\end{equation}
\begin{equation}
  w_{dskmpj}\le y_{kmpj}, \qquad \forall d\in R,\ s\in S,\ k\in K,\ m\in H,\ p\in P,\ j\in N.
\end{equation}
\begin{equation}
  \sum_{k\in K}\sum_{m\in H}\sum_{p\in P}\sum_{j\in N}w_{dskmpj} =\eta_{ds}, \qquad \forall d\in R,\ s\in S.
\end{equation}
\begin{equation}
  \eta_{d,s+1}\le \eta_{ds}, \qquad \forall d\in R,\ s=1,\ldots,S_{\max}-1.
\end{equation}

\subsubsection{Truck capacity và tải chuyển tiếp giữa các chuyến}

\begin{equation}
  L_{k10}=\sum_{i\in N^D}q_i a_{k1i}, \qquad \forall k\in K.
\end{equation}
Từ chuyến thứ hai trở đi, tải đầu chuyến gồm hàng pickup C2 còn trên truck từ chuyến
trước cộng với hàng C1 mới được chất ở depot:
\begin{equation}
  L_{km0}=L^{end}_{k,m-1}+\sum_{i\in N^D}q_i a_{kmi}, \qquad \forall k\in K,\ m=2,\ldots,H_{\max}.
\end{equation}
\begin{equation}
  L_{km0}\le Q_Tg_{km}, \qquad \forall k\in K,\ m\in H.
\end{equation}
Tải truck sau khi hoàn tất phục vụ và bàn giao tại vị trí $p$:
\begin{equation}
  L_{kmp} = L_{km,p-1} +\sum_{i\in N^P}q_i y_{kmpi} -\sum_{i\in N^{DL}}q_i y_{kmpi} -\sum_{i\in N^D}q_i y_{kmpi} +\sum_{i\in N^D}q_i b_{kmpi},
\end{equation}
\[
  \forall k\in K,\ m\in H,\ p\in P.
\]
\begin{equation}
  0\le L_{kmp}\le Q_Tu_{kmp}, \qquad \forall k\in K,\ m\in H,\ p\in P.
\end{equation}
Tải cuối chuyến bằng tải tại khách cuối cùng của chuyến:
\begin{equation}
  L^{end}_{km}\ge L_{kmp}-M(1-\lambda_{kmpi}), \qquad \forall k\in K,\ m\in H,\ p\in P,\ i\in N.
\end{equation}
\begin{equation}
  L^{end}_{km}\le L_{kmp}+M(1-\lambda_{kmpi}), \qquad \forall k\in K,\ m\in H,\ p\in P,\ i\in N.
\end{equation}
\begin{equation}
  0\le L^{end}_{km}\le Q_Tg_{km}, \qquad \forall k\in K,\ m\in H.
\end{equation}
Để cấm tường minh hàng C1 đi qua depot nhưng vẫn cho phép hàng C2 đi qua depot, tải
cuối chuyến phải bằng đúng cân bằng pickup--delivery C2 lũy kế:
\begin{equation}
  L^{end}_{km} = \sum_{(p',\ell)\in\mathcal A}q_{p'}\left(\sum_{r=1}^{m}\sum_{s\in P}y_{krs p'} - \sum_{r=1}^{m}\sum_{s\in P}y_{krs\ell}\right), \qquad \forall k\in K,\ m\in H.
\end{equation}
Vì pickup và delivery cùng truck và pickup đứng trước delivery, vế phải không âm.
Nếu $L^{end}_{km}>0$, toàn bộ phần tải đó là hàng C2 đã pickup nhưng chưa tới
delivery tương ứng. Do đó C2 có thể được giữ qua depot, còn C1 thì không.

\subsubsection{Thời gian truck}

\begin{equation}
  y_{km1i}=1\ \Longrightarrow\ A_{km1}=T^0_{km}+t^T_{0i}, \qquad \forall k\in K,\ m\in H,\ i\in N.
\end{equation}
\begin{equation}
  z_{kmpij}=1\ \Longrightarrow\ A_{km,p+1}=D_{kmp}+t^T_{ij}, \qquad \forall k\in K,\ m\in H,\ p=1,\ldots,P_{\max}-1,\ i,j\in N.
\end{equation}
\begin{equation}
  D_{kmp}\ge A_{kmp}+s_Tu_{kmp}, \qquad \forall k\in K,\ m\in H,\ p\in P.
\end{equation}
Khi không có drone resupply tại vị trí đang dùng, thời điểm rời được khóa đúng bằng
thời điểm đến cộng service time:
\begin{equation}
  D_{kmp}\le A_{kmp}+s_T+M\rho_{kmp}+M(1-u_{kmp}), \qquad \forall k\in K,\ m\in H,\ p\in P.
\end{equation}
\begin{equation}
  D_{kmp}\ge F^{arr}_{ds}+\delta_m -M\left(1-\sum_{j\in N}w_{dskmpj}\right), \qquad \forall d\in R,\ s\in S,\ k\in K,\ m\in H,\ p\in P.
\end{equation}
\begin{equation}
  \lambda_{kmpi}=1\ \Longrightarrow\ B_{km}=D_{kmp}+t^T_{i0}, \qquad \forall k\in K,\ m\in H,\ p\in P,\ i\in N.
\end{equation}
Chuyến đầu tiên không cộng thời gian bốc dỡ truck tại depot:
\begin{equation}
  T^0_{k1}\ge r_i a_{k1i}, \qquad \forall k\in K,\ i\in N^D.
\end{equation}
Từ chuyến thứ hai, truck chỉ rời depot sau khi chuyến trước kết thúc và đã bốc dỡ:
\begin{equation}
  T^0_{km}\ge B_{k,m-1}+\delta_T-M(1-g_{km}), \qquad \forall k\in K,\ m=2,\ldots,H_{\max}.
\end{equation}
Release date cho hàng C1 chất lên truck tại depot đầu chuyến:
\begin{equation}
  T^0_{km}\ge r_i a_{kmi}, \qquad \forall k\in K,\ m\in H,\ i\in N^D.
\end{equation}

\subsubsection{Release date, drone và đồng bộ}

\begin{equation}
  F^{dep}_{ds}\ge r_i+\delta_d -M\left(1-b_{kmpi}\right) -M\left(1-\sum_{j\in N}w_{dskmpj}\right), \qquad \forall d\in R,\ s\in S,\ k\in K,\ m\in H,\ p\in P,\ i\in N^D.
\end{equation}
\begin{equation}
  F^{dep}_{d,s+1}\ge F^{ret}_{ds}+\delta_d-M(1-\eta_{d,s+1}), \qquad \forall d\in R,\ s=1,\ldots,S_{\max}-1.
\end{equation}
\begin{equation}
  w_{dskmpj}=1\ \Longrightarrow\ F^{arr}_{ds}=F^{dep}_{ds}+t^D_{0j}, \qquad \forall d\in R,\ s\in S,\ k\in K,\ m\in H,\ p\in P,\ j\in N.
\end{equation}
\begin{equation}
  w_{dskmpj}=1\ \Longrightarrow\ F^{ret}_{ds}=F^{arr}_{ds}+\delta_m+t^D_{j0}, \qquad \forall d\in R,\ s\in S,\ k\in K,\ m\in H,\ p\in P,\ j\in N.
\end{equation}
Drone không cần chờ truck tại điểm hẹn. Vì vậy nếu chuyến bay $(d,s)$ được gán
cho điểm hẹn $(k,m,p)$, thời điểm drone tới điểm hẹn không được sớm hơn thời điểm
truck tới vị trí đó:
\begin{equation}
  F^{arr}_{ds}\ge A_{kmp}-M\left(1-\sum_{j\in N}w_{dskmpj}\right), \qquad \forall d\in R,\ s\in S,\ k\in K,\ m\in H,\ p\in P.
\end{equation}
Truck có thể chờ drone tại điểm hẹn, nên mô hình không dùng tham số giới hạn thời
gian chờ; nghiệm tối ưu sẽ chọn $F^{dep}_{ds}$ phù hợp với release date, thứ tự
chuyến bay của drone, và thời điểm truck tới điểm hẹn.
\begin{equation}
  F^{ret}_{ds}-F^{dep}_{ds}\le E+M(1-\eta_{ds}), \qquad \forall d\in R,\ s\in S.
\end{equation}

\subsubsection{$C_{\max}$}

\begin{equation}
  C_{\max}\ge B_{km}, \qquad \forall k\in K,\ m\in H.
\end{equation}
\begin{equation}
  C_{\max}\ge F^{ret}_{ds}, \qquad \forall d\in R,\ s\in S.
\end{equation}

\subsubsection{Miền giá trị}

\begin{align*}
y_{kmpi},u_{kmp},g_{km},z_{kmpij},\lambda_{kmpi}, a_{kmi},b_{kmpi},\rho_{kmp},w_{dskmpj},\eta_{ds} &\in\{0,1\},\\
T^0_{km},A_{kmp},D_{kmp},B_{km}, F^{dep}_{ds},F^{arr}_{ds},F^{ret}_{ds},L_{kmp},L^{end}_{km},C_{\max} &\ge0.
\end{align*}

\subsection{Ghi chú triển khai}

Mô hình trên và file \texttt{MILP.py} cùng sử dụng truck nhiều chuyến, truck và drone
có capacity, C1 phải được nhận và giao trong cùng chuyến truck, hàng pickup C2 có thể
được giữ qua depot để giao ở chuyến sau, và mỗi chuyến drone chỉ có một rendezvous.
Khi CPLEX dừng ở time limit, incumbent có thể chứa slack trong các biến thời gian dù
quyết định route/load/resupply đã cố định. Vì vậy code cố định toàn bộ biến nhị phân
của incumbent và giải lại bài toán thời gian liên tục trước khi báo objective hoặc so
sánh với decoder heuristic.

%==============================================================================
\section{Thuật toán Giải}
\label{sec:algorithm}
%==============================================================================

\subsection{Tổng quan: kiến trúc hai lớp TS-ALNS-LSP}

Thuật toán heuristic hiện tại (\texttt{pdp\_tabu\_main.cpp}, \texttt{pdp\_lsp.cpp},
\texttt{pdp\_alns.cpp}, \texttt{pdp\_lb.cpp}) \emph{không phải} GA+Tabu như mô tả
trong các phiên bản báo cáo trước, mà là một kiến trúc \textbf{Tabu Search thích nghi
(Adaptive Tabu Search -- ATS)} hai lớp, tách biệt quyết định \emph{thứ tự phục vụ
khách hàng} khỏi quyết định \emph{nguồn hàng/resupply}:
\begin{itemize}
    \item \textbf{Lớp 1 -- Routing (Tabu Search + ALNS escape):} làm việc trên
    \texttt{VRPSolution} -- chỉ thứ tự khách hàng của mỗi truck, không chứa thông tin
    loading. Không gian tìm kiếm là hoán vị khách hàng.
    \item \textbf{Lớp 2 -- Loading Subproblem (LSP, CP-SAT):} với route truck đã cố
    định ở Lớp 1, giải một bài toán con bằng CP-SAT (OR-Tools) để quyết định mỗi
    kiện C1 được chất ở depot, hay được drone resupply tại vị trí nào, nhằm
    tối thiểu $C_{\max}$ cho route đó.
\end{itemize}
Việc tách lớp này cho phép Lớp 1 dùng một \textbf{cận dưới (lower bound)} rẻ để
sàng lọc hầu hết các move trước khi gọi Lớp 2 (tốn kém hơn nhiều vì phải giải CP-SAT).

\subsection{Cận dưới (Lower Bound) cho một route}

Với route $\chi_v = (c_1, \ldots, c_m)$ của truck $v$, cận dưới thời điểm truck đến
khách $c_k$ được tính truyền tiến (\texttt{pdp\_lb.cpp}):
\begin{equation}
  \ell_i(c) = r_c + \min\big(\delta_d + t^D_{0c},\; t^T_{0c}\big) + s_T
  \qquad\text{(cận dưới tại khách đầu tiên, có thể nhận hàng bằng drone hoặc truck)}
\end{equation}
\begin{equation}
  \ell_i(c_k) = \max\big(\ell_i(c_k),\; \ell_i(c_{k-1}) + t^T_{c_{k-1}c_k}\big),
  \qquad k = 2, \ldots, m,
\end{equation}
\begin{equation}
  LB(\chi_v) = \ell_i(c_m) + t^T_{c_m 0}.
\end{equation}
Cận dưới toàn nghiệm là $LB(\chi) = \max_v LB(\chi_v)$. Cận dưới này được dùng
để: (i) chọn truck "bottleneck" (route có LB lớn nhất) làm ưu tiên tìm move; (ii)
tính ngưỡng chấp nhận move; (iii) làm \texttt{ub\_cutoff} để CP-SAT dừng sớm khi
không thể cải thiện nghiệm tốt nhất.

\subsection{Lớp 1: Tabu Search thích nghi}

\begin{algorithm}[H]
\caption{Adaptive Tabu Search (Lớp 1 -- routing)}
\begin{algorithmic}[1]
\STATE \textbf{Input:} $\chi_0$ (VRP khởi tạo), cấu hình \texttt{TSConfig}
\STATE $\chi^\star \leftarrow \chi_0$;\quad $\chi \leftarrow \chi_0$
\STATE $(z^\star, LSP^\star) \leftarrow$ \textsc{EvaluateRoutes}$(\chi^\star)$ \COMMENT{gọi Lớp 2}
\STATE $LB^\star \leftarrow LB(\chi^\star)$;\quad $tabu \leftarrow \emptyset$;\quad $w_{\text{op}} \leftarrow (1,1,1)$ \COMMENT{RELOCATE, SWAP, 2-OPT}
\WHILE{chưa hết \texttt{time\_limit\_sec}}
    \STATE $\theta \sim \mathcal{U}(\theta_{lo}, \theta_{hi})$ \COMMENT{tenure ngẫu nhiên hoá, $\theta_{base}=\mathrm{round}(5\sqrt{n})$ [Taillard 1993]}
    \STATE $\varphi \leftarrow \min(1, \text{stagnation}/\eta_{\max})$;\quad $\tau \leftarrow (1-\varphi)z^\star + \varphi\, LB^\star$
    \STATE chọn toán tử $op$ theo roulette-wheel trên $w_{\text{op}}$
    \STATE tìm move tốt nhất của $op$ trên truck "bottleneck" (LB lớn nhất) trước, rồi các truck còn lại
    \STATE \textbf{nếu} không tìm được move, thử lần lượt 2 toán tử còn lại (fallback)
    \STATE move được nhận nếu $LB(\chi') < z^\star$ (aspiration) hoặc ($\neg$tabu và $LB(\chi') < \tau$)
    \IF{tìm được move $m^\star$}
        \STATE $\chi \leftarrow \chi \oplus m^\star$;\quad $tabu[m^\star] \leftarrow iter+\theta$
        \STATE $(z, LSP) \leftarrow$ \textsc{EvaluateRoutes}$(\chi,\ \text{cutoff}=z^\star-\varepsilon)$ \COMMENT{chỉ chấp nhận nếu $< z^\star-\varepsilon$}
        \IF{$LSP$ khả thi và $z < z^\star$}
            \STATE $z^\star \leftarrow z$;\ $\chi^\star \leftarrow \chi$;\ $LB^\star \leftarrow LB(\chi^\star)$;\ stagnation $\leftarrow 0$;\ score $\leftarrow 3$
        \ELSE
            \STATE stagnation++;\ score $\leftarrow 1$
        \ENDIF
        \STATE cập nhật $w_{\text{op}}[op] \leftarrow \max(0.05,\ (1-\rho)w_{\text{op}}[op] + \rho\cdot\text{score})$ \COMMENT{$\rho=0.1$}
    \ELSE
        \STATE stagnation++;\ cập nhật $w_{\text{op}}[op]$ với score $=0$
    \ENDIF
    \IF{stagnation $\ge \eta_{\max}$ (mặc định 150)}
        \STATE gọi \textsc{ALNSEscape}$(\chi^\star)$ (xem Thuật toán~\ref{alg:alns}); $tabu \leftarrow \emptyset$;\ stagnation $\leftarrow 0$
    \ENDIF
    \IF{không cải thiện toàn cục trong 500 iteration liên tiếp}
        \STATE ép stagnation $=\eta_{\max}$ để kích hoạt ALNS escape ngay (restart) [Sacramento et al., 2019]
    \ENDIF
\ENDWHILE
\RETURN $\chi^\star, z^\star, LSP^\star$
\end{algorithmic}
\end{algorithm}

\textbf{3 toán tử lân cận (neighbourhood moves):} RELOCATE (chuyển 1 khách sang vị
trí/truck khác), SWAP (hoán đổi 2 khách, có thể khác truck), 2-OPT (đảo một đoạn
trong cùng route). Khách loại \texttt{DL} (delivery của cặp C2) không bao giờ bị
RELOCATE/SWAP trực tiếp -- di chuyển node \texttt{P} tương ứng sẽ kéo theo vị trí hợp
lệ của \texttt{DL}; mọi move được kiểm tra lại ràng buộc tiền tố C2 (pickup trước
delivery, cùng truck) trước khi được xem xét.

\subsection{Lớp 2: Loading Subproblem (CP-SAT)}

Với route cố định, \texttt{pdp\_lsp.cpp} xây dựng một mô hình CP-SAT nhỏ hơn nhiều
so với MILP đầy đủ (không còn biến định tuyến $y, z$), chỉ còn quyết định
\textbf{3 cơ chế nạp hàng} cho mỗi kiện C1:
\begin{enumerate}
    \item \textbf{Slot 0} -- chất tại depot đầu chuyến ($a_{kmi}$ trong MILP);
    \item \textbf{RS slot} tại vị trí $(v,k)$ -- drone resupply tại vị trí đó
    ($b_{kmpi}$ trong MILP);
    \item \textbf{DR slot} tại vị trí $(v,k)$, $k\ge 1$ -- truck quay lại depot giữa
    chừng trước vị trí đó để nhận hàng.
\end{enumerate}
Vì truck trong lớp này được coi là \emph{không giới hạn tải} (giả định đơn giản hoá
của tầng heuristic, khác với MILP đầy đủ có $Q_T$), ràng buộc chính chỉ là release
date: một kiện chỉ được nạp (ở bất kỳ cơ chế nào) sau thời điểm $r_i$. Quy trình giải:
\begin{enumerate}
    \item \texttt{greedyLoading}: chất toàn bộ C1 tại depot (luôn khả thi), dùng làm
    \emph{warm start} (\texttt{AddHint}) cho CP-SAT;
    \item \texttt{solveLSP}: CP-SAT tìm cách gán lại một số kiện sang RS/DR để giảm
    $C_{\max}$, với \texttt{cutoff\_sec} (mặc định 3s/lần gọi) và \texttt{ub\_cutoff}
    = $z^\star - \varepsilon$ để dừng sớm nếu không thể cải thiện nghiệm tốt nhất
    hiện tại;
    \item \texttt{evaluateRoutes} = \texttt{greedyLoading} + \texttt{solveLSP}, trả
    về \texttt{LSPResult} (khả thi hay không, $C_{\max}$, và \texttt{SolutionEncoding}
    đầy đủ nếu giải được).
\end{enumerate}

\subsection{ALNS escape (thoát cục bộ)}

\begin{algorithm}[H]
\caption{ALNS escape (gọi khi Tabu Search trì trệ)}\label{alg:alns}
\begin{algorithmic}[1]
\STATE $n_{remove} \sim \mathcal{U}(0.15n,\ 0.35n)$ [Ropke \& Pisinger, 2006]
\FOR{$t = 1$ đến $n_{\text{alns\_escape}}$ (mặc định 5)}
    \STATE chọn toán tử phá (destroy) $d$ và sửa (repair) $r$ theo roulette-wheel
    \STATE $\chi' \leftarrow r(d(\chi^\star, n_{remove}))$
    \STATE \textbf{nếu} $\chi'$ vi phạm tiền tố C2 thì bỏ qua lần thử này
    \STATE $(z', LSP') \leftarrow$ \textsc{EvaluateRoutes}$(\chi',\ \text{cutoff}=z^\star-\varepsilon)$
    \IF{$LSP'$ khả thi và $z' < z^\star$}
        \STATE $z^\star \leftarrow z'$;\ $\chi^\star \leftarrow \chi'$;\ score $\leftarrow 2$ (new best)
    \ELSIF{$LSP'$ khả thi và $z' < 1.05\cdot z^\star$}
        \STATE score $\leftarrow 1$ (chấp nhận nhẹ)
    \ELSE
        \STATE score $\leftarrow 0$
    \ENDIF
    \STATE cập nhật trọng số $d, r$ theo $w \leftarrow (1-\rho)w + \rho\cdot\{0.5,1.0,2.0\}[\text{score}]$
\ENDFOR
\end{algorithmic}
\end{algorithm}

\textbf{3 toán tử phá (destroy):}
\begin{enumerate}
    \item \textit{Random segment}: xoá một đoạn liên tiếp ngẫu nhiên trên 1 truck;
    \item \textit{Worst removal}: xoá các khách có "chi phí lệch tuyến" (detour cost
    $t(prev,n){+}t(n,next){-}t(prev,next)$) lớn nhất, có thêm nhiễu ngẫu nhiên trong
    top $3\times n_{remove}$ để tránh luôn chọn đúng 1 tập;
    \item \textit{Related removal}: chọn 1 khách ngẫu nhiên làm "hạt giống", xoá thêm
    các khách gần nhất về khoảng cách (Shaw removal đơn giản hoá).
\end{enumerate}
\textbf{2 toán tử sửa (repair):} \textit{Greedy insert} (chèn từng khách bị xoá vào
vị trí rẻ nhất, xử lý C1 trước rồi C2) và \textit{Regret-2 insert} (ưu tiên chèn
khách có "regret" lớn nhất -- chênh lệch giữa vị trí chèn tốt nhất và tốt nhì -- để
tránh kẹt các khách khó chèn về sau). Cả hai đều tôn trọng ràng buộc C2 (P phải đã có
trong route trước khi chèn DL tương ứng).

\subsection{Tổng kết tham số mặc định (\texttt{TSConfig})}

\begin{table}[H]
\centering
\begin{tabular}{ll}
\toprule
Tham số & Giá trị mặc định / ý nghĩa\\
\midrule
\texttt{eta\_max} & 150 -- số iteration trì trệ trước khi kích hoạt ALNS escape\\
\texttt{no\_improve\_limit} & 500 -- số iteration không cải thiện toàn cục trước khi ép restart\\
\texttt{n\_alns\_escape} & 5 -- số lần thử destroy/repair mỗi lần escape\\
\texttt{cutoff\_eps} ($\varepsilon$) & 0.5 -- chỉ chấp nhận LSP nếu $z < z^\star-\varepsilon$\\
\texttt{lsp\_cutoff} & 3.0s -- thời gian CP-SAT tối đa mỗi lần gọi\\
\texttt{time\_limit\_sec} & ngân sách tổng theo wall-clock (2.5 phút bộ nhỏ, 60 phút bộ lớn)\\
\bottomrule
\end{tabular}
\end{table}

\emph{Ghi chú:} kiến trúc và các tham số trên bám theo khung Adaptive Tabu Search
được đề xuất trong Tran et al. (2026) -- xem mục~\ref{sec:related-work}. Tuy nhiên,
khác với mô hình VRP-MRDR \emph{multi-rendezvous} của bài báo đó, Lớp 2 hiện tại
(\texttt{pdp\_lsp.cpp}) chỉ gán \textbf{một} điểm hẹn cho mỗi chuyến bay drone (RS
slot), khớp với mô hình \emph{one-rendezvous} ở mục~2 chứ chưa hỗ trợ một drone tiếp
tế tuần tự nhiều truck trong cùng một chuyến bay.

%==============================================================================
\section{Hàm Decode - Chuyển đổi Sequence thành Solution}
%==============================================================================

\subsection{Tổng quan}

Hàm decode là thành phần quan trọng nhất, chuyển đổi một chromosome (sequence) thành một solution hoàn chỉnh với:
\begin{itemize}
    \item Routes cho từng xe tải
    \item Timeline chi tiết (arrival time, departure time)
    \item Các drone resupply events
    \item Tính toán C\_max chính xác
\end{itemize}

\subsection{Xử lý các loại khách hàng}

\subsubsection{Khách hàng loại P (Pickup)}

\begin{itemize}
    \item Xe tải đến điểm pickup, load hàng lên xe
    \item Đánh dấu cặp pickup-delivery đã được pickup
    \item Hàng được lưu trên xe cho đến khi delivery
\end{itemize}

\subsubsection{Khách hàng loại DL (Delivery Location)}

\begin{itemize}
    \item Tìm xe đã pickup cặp tương ứng
    \item Xe đó BẮT BUỘC phải giao tại DL này
    \item Không cần đợi ready time (hàng đã có trên xe)
\end{itemize}

\subsubsection{Khách hàng loại D (Drone-eligible)}

Có 2 phương án được so sánh:

\textbf{Phương án DJ1 - Drone Resupply với Consolidation:}
\begin{enumerate}
    \item Tìm các customers có thể gom chung (consolidation candidates)
    \item Đánh giá hiệu quả của từng batch size
    \item Chọn phương án có effective cost thấp nhất
\end{enumerate}

\textbf{Phương án DJ2 - Truck giao trực tiếp:}
\begin{enumerate}
    \item Truck về depot
    \item Đợi ready time của khách hàng
    \item Lấy hàng và đi giao
\end{enumerate}

\subsection{Consolidation - Gom nhiều khách hàng}

Điều kiện để consolidate nhiều khách hàng type D:

\begin{equation}
    |r_i - r_j| \leq 30 \text{ phút} \quad \land \quad d_{ij}^{drone} \leq 10 \text{ km}
\end{equation}

Trong đó:
\begin{itemize}
    \item $r_i, r_j$: Ready time của hai khách hàng
    \item $d_{ij}^{drone}$: Khoảng cách Euclidean giữa hai khách hàng
\end{itemize}

\subsection{Effective Cost - Đánh giá lợi ích thực sự}

Thay vì chỉ so sánh completion time, chương trình sử dụng effective cost:

\begin{equation}
    effective\_cost = completion\_time - savings
\end{equation}

\begin{equation}
    savings = (n - 1) \times trip\_overhead
\end{equation}

Trong đó:
\begin{itemize}
    \item $n$: Số khách hàng trong batch
    \item $trip\_overhead \approx 15$ phút (depot load time + thời gian bay về)
\end{itemize}

\textbf{Ví dụ:}
\begin{itemize}
    \item Giao 2 khách riêng lẻ: $45 + 30 = 75$ phút (effective)
    \item Giao 2 khách cùng lúc: $50 - 15 = 35$ phút (effective)
    \item $\Rightarrow$ Gom nhóm tiết kiệm 40 phút!
\end{itemize}

\subsection{Smart Merge}

Khi merge 2 drone trips, có kiểm tra thông minh:

\begin{equation}
    \text{Merge allowed} \Leftrightarrow \neg(ready\_gap > 15 \land avg\_dist > 3)
\end{equation}

Ma trận quyết định:
\begin{table}[H]
\centering
\caption{Quyết định merge dựa trên ready gap và distance}
\begin{tabular}{ccc}
\toprule
\textbf{Ready Gap} & \textbf{Avg Distance} & \textbf{Quyết định} \\
\midrule
$\leq 15$ phút & Bất kỳ & Có thể merge \\
$> 15$ phút & $\leq 3$ km & Có thể merge \\
$> 15$ phút & $> 3$ km & KHÔNG merge \\
\bottomrule
\end{tabular}
\end{table}

\subsection{Propagate Delay}

Đây là bước quan trọng để tính đúng C\_max:

\begin{algorithm}[H]
\caption{Propagate Delay}
\begin{algorithmic}[1]
\FOR{each resupply event}
    \STATE Find resupply point in truck route
    \STATE $required\_departure \leftarrow resupply\_end\_time$
    \IF{$required\_departure > truck.departure[k]$}
        \STATE $delay \leftarrow required\_departure - truck.departure[k]$
        \STATE $truck.departure[k] \leftarrow required\_departure$
        \FOR{all subsequent nodes $m > k$}
            \STATE $truck.arrival[m] \leftarrow truck.arrival[m] + delay$
            \STATE $truck.departure[m] \leftarrow truck.departure[m] + delay$
        \ENDFOR
    \ENDIF
\ENDFOR
\end{algorithmic}
\end{algorithm}

%==============================================================================
\section{Cấu trúc Chương trình}
%==============================================================================

\subsection{Tổ chức mã nguồn}

\begin{table}[H]
\centering
\caption{Các file mã nguồn chính}
\begin{tabular}{ll}
\toprule
\textbf{File} & \textbf{Chức năng} \\
\midrule
\texttt{pdp\_types.h} & Định nghĩa cấu trúc dữ liệu \\
\texttt{pdp\_reader.cpp/h} & Đọc dữ liệu từ file instance \\
\texttt{pdp\_init.cpp/h} & Khởi tạo quần thể \\
\texttt{pdp\_fitness.cpp/h} & Hàm decode và tính fitness \\
\texttt{pdp\_ga.cpp/h} & Các toán tử GA \\
\texttt{pdp\_tabu.cpp/h} & Thuật toán Tabu Search \\
\texttt{pdp\_localsearch.cpp/h} & Các operator local search \\
\texttt{pdp\_validation.cpp/h} & Kiểm tra tính hợp lệ của solution \\
\texttt{main\_ga\_tabu.cpp} & Chương trình chính \\
\bottomrule
\end{tabular}
\end{table}

\subsection{Các cấu trúc dữ liệu chính}

\subsubsection{PDPData - Dữ liệu bài toán}

\begin{lstlisting}[language=C++]
struct PDPData {
    int numNodes, numCustomers;
    int depotIndex;
    vector<double> xCoords, yCoords;
    vector<string> nodeTypes;    // "DEPOT", "P", "DL", "D"
    vector<int> readyTimes;
    vector<int> pairIds;
    vector<int> demands;
    
    int numTrucks = 2;
    int numDrones = 2;
    double droneSpeed = 60.0;
    double truckSpeed = 40.0;
    double droneEndurance = 90.0;
    
    vector<vector<double>> truckDistMatrix;  // Manhattan
    vector<vector<double>> droneDistMatrix;  // Euclidean
};
\end{lstlisting}

\subsubsection{PDPSolution - Lời giải}

\begin{lstlisting}[language=C++]
struct PDPSolution {
    double totalCost;           // C_max
    double totalPenalty;
    bool isFeasible;
    
    vector<vector<int>> routes;
    vector<TruckRouteInfo> truck_details;
    vector<ResupplyEvent> resupply_events;
    vector<double> drone_completion_times;
};
\end{lstlisting}

\subsubsection{ResupplyEvent - Sự kiện tiếp tế}

\begin{lstlisting}[language=C++]
struct ResupplyEvent {
    vector<int> customer_ids;   // Customers in this trip
    int drone_id;
    int truck_id;
    int resupply_point;
    
    double drone_depart_time;
    double drone_arrive_time;
    double truck_arrive_time;
    double resupply_start_time;
    double resupply_end_time;
    double drone_return_time;
    double truck_delivery_end;
    double total_flight_time;
};
\end{lstlisting}

%==============================================================================
\section{Validation - Kiểm tra Tính hợp lệ}
%==============================================================================

Chương trình thực hiện 6 bước kiểm tra để đảm bảo solution hợp lệ:

\subsection{Customer Coverage}
Kiểm tra tất cả khách hàng đều được phục vụ:
\begin{equation}
    |served\_customers| = |C_P \cup C_{DL} \cup C_D|
\end{equation}

\subsection{Drone Endurance}
Với mỗi drone trip:
\begin{equation}
    total\_flight\_time \leq E = 90 \text{ phút}
\end{equation}

\subsection{Timeline Consistency}
Với mỗi xe tải, timeline phải nhất quán:
\begin{equation}
    arrival[i] \leq departure[i] \leq arrival[i+1]
\end{equation}

\subsection{Drone-Truck Synchronization}
Tại điểm resupply:
\begin{equation}
    resupply\_start = \max(drone\_arrive, truck\_arrive)
\end{equation}

\subsection{C\_max Verification}
\begin{equation}
    C_{max}^{reported} = C_{max}^{calculated}
\end{equation}

\subsection{Depot Returns}
Kiểm tra hiệu quả sử dụng drone (informational).

%==============================================================================
\section{Kết quả Thực nghiệm}
%==============================================================================

\subsection{Môi trường thực nghiệm}

\begin{itemize}
    \item \textbf{Ngôn ngữ}: C++17
    \item \textbf{Compiler}: clang++ với -O2 optimization
    \item \textbf{Hệ điều hành}: macOS
    \item \textbf{Parameters}: Population = 200, Generations = 100
\end{itemize}

\subsection{Kết quả trên bộ dữ liệu U\_30}

\begin{table}[H]
\centering
\caption{Kết quả thực nghiệm trên instances U\_30\_0.5}
\begin{tabular}{lccccc}
\toprule
\textbf{Instance} & \textbf{Customers} & \textbf{C\_max (phút)} & \textbf{Drone Trips} & \textbf{Valid} \\
\midrule
U\_30\_0.5\_Num\_1 & 30 & 315.81 & 16 & \checkmark \\
U\_30\_0.5\_Num\_2 & 29 & 287.12 & 13 & \checkmark \\
U\_30\_0.5\_Num\_3 & 30 & 292.02 & 13 & \checkmark \\
\bottomrule
\end{tabular}
\end{table}

\subsection{Phân tích Consolidation}

Consolidation giúp giảm đáng kể số chuyến drone:

\begin{table}[H]
\centering
\caption{Phân tích hiệu quả consolidation}
\begin{tabular}{lcccc}
\toprule
\textbf{Instance} & \textbf{Trips 1 cust} & \textbf{Trips 2 custs} & \textbf{Trips 3+ custs} & \textbf{Tổng} \\
\midrule
Num\_1 & 7 & 7 & 2 & 16 \\
Num\_2 & 4 & 4 & 5 & 13 \\
Num\_3 & 3 & 5 & 5 & 13 \\
\bottomrule
\end{tabular}
\end{table}

\subsection{Nhận xét}

\begin{enumerate}
    \item \textbf{Tính khả thi}: 100\% solutions đều valid với C\_max reported = calculated
    \item \textbf{Consolidation hiệu quả}: Trung bình 60-70\% drone trips gom được 2+ customers
    \item \textbf{Drone endurance}: Tất cả trips đều trong giới hạn 90 phút
    \item \textbf{Sync chính xác}: Drone-truck synchronization hoạt động đúng
\end{enumerate}

%==============================================================================
\section{Kết luận}
%==============================================================================

\subsection{Đóng góp chính}

Chương trình đã đạt được các mục tiêu đề ra:

\begin{enumerate}
    \item \textbf{Thuật toán lai ghép hiệu quả}: Kết hợp GA với Tabu Search, sử dụng adaptive operators để cân bằng exploration và exploitation
    
    \item \textbf{Mô hình One-Rendezvous với Consolidation}: Cho phép drone giao nhiều packages trong một chuyến, giảm đáng kể số trips cần thiết
    
    \item \textbf{Effective Cost và Smart Merge}: Đánh giá lợi ích thực sự của consolidation, tránh merge không hiệu quả
    
    \item \textbf{Propagate Delay}: Tính toán chính xác C\_max với đồng bộ hóa drone-truck
    
    \item \textbf{Validation toàn diện}: 6 bước kiểm tra đảm bảo solution hợp lệ
\end{enumerate}

\subsection{Hướng phát triển}
\label{sec:future-work}

\begin{enumerate}
    \item Mở rộng sang mô hình Multiple-Rendezvous (VRP-MRDR, Tran et al. 2026) -- cho
    phép một drone tiếp tế tuần tự nhiều truck trong cùng một chuyến bay, thay vì chỉ
    một điểm hẹn như hiện tại
    \item Tích hợp Machine Learning để dự đoán operator hiệu quả
    \item Parallelization cho large-scale instances
    \item Thêm các ràng buộc thực tế: time windows, traffic, weather
\end{enumerate}

%==============================================================================
\section*{Tài liệu tham khảo}
%==============================================================================

\begin{enumerate}
    \item Pina-Pardo, J.C., Silva, F., Smith, A.E. (2021). The traveling salesman
    problem with release dates and drone resupply. \emph{Computers \& Operations
    Research}, 129, 105170.
    \item Pina-Pardo, J.C., Silva, F., Smith, A.E., Gatica, R.A. (2024a). Dynamic
    vehicle routing problem with drone resupply. \emph{Transportation Research Part
    C: Emerging Technologies}, 162, 104611.
    \item Pina-Pardo, J.C., Silva, F., Smith, A.E., Gatica, R.A. (2024b). Fleet
    resupply by drones for last-mile delivery. \emph{European Journal of Operational
    Research}, 316(1), 168--182.
    \item Tran, T.H., Nguyen, P.K., Binh, H.T.T., Tran, H.H., Nguyen, S. (2026).
    Coordinated truck routing and multi-rendezvous drone resupply under shipment
    release dates. \emph{Transportation Research Part E: Logistics and
    Transportation Review}, 215, 105082.
    \item Taillard, E. (1993). Parallel iterative search methods for vehicle routing
    problems. \emph{Networks}, 23(8), 661--673.
    \item Ropke, S., Pisinger, D. (2006). An adaptive large neighborhood search
    heuristic for the pickup and delivery problem with time windows.
    \emph{Transportation Science}, 40(4), 455--472.
    \item Sacramento, D., Pisinger, D., Ropke, S. (2019). An adaptive large
    neighborhood search metaheuristic for the vehicle routing problem with drones.
    \emph{Transportation Research Part C: Emerging Technologies}, 102, 289--315.
\end{enumerate}

\end{document}
